% Chapter 4: PVE-ZKP Architecture

\chapter{PVE-ZKP Architecture}

This chapter describes the architecture of PVE-ZKP in terms of components,
interfaces, and message flows. It mirrors the concise architecture presented in
the conference paper, but provides additional detail and discussion of design
choices and deployment variants.

\section{High-Level Overview}

At a high level, PVE-ZKP interposes a validation service between the browser
client and the receiving system. The validation service runs zero-knowledge
policy checks on raw form data and returns only an opaque proof identifier to
the browser. The receiving system uses this identifier to fetch and verify
proofs before deciding whether to accept the submission.

The architecture can be summarized textually as:

\begin{quote}
Browser (client) $\rightarrow$ \texttt{validator.js}:3001 (proof generation,
proof storage under \texttt{proofId}) $\rightarrow$ browser (forward
\texttt{proofId}) $\rightarrow$ \texttt{receiving\_system.js}:3002 (fetch and
verify proofs) $\rightarrow$ accept / reject.
\end{quote}

In the prototype implementation, both the validation service and the receiving
system are implemented as Node.js applications, but the architectural pattern
is independent of the specific technology stack.

\section{Component Responsibilities}

\subsection{Browser Client}

The browser client is responsible for collecting form inputs from the user and
communicating with both services:

\begin{itemize}[leftmargin=*]
  \item \textbf{To validator:} send form data as JSON to the validation service
    at \texttt{/validate} on port 3001 and handle the returned
    \texttt{proofId}.
  \item \textbf{To receiver:} send a follow-up request containing only
    \texttt{proofId} to the receiving system at \texttt{/receive} on port 3002.
\end{itemize}

The client performs basic sanity checks for user experience, but security does
not rely on client-side validation.

\subsection{Validation Service}

The validation service, implemented in \texttt{validator.js}, exposes an HTTP
API for proving and for proof retrieval. It:

\begin{itemize}[leftmargin=*]
  \item accepts raw form data from the browser;
  \item maps fields to inputs of the five Circom policy circuits;
  \item generates Groth16 proofs using snarkjs;
  \item stores proofs in an in-memory map keyed by a randomly generated
        \texttt{proofId}; and
  \item returns \texttt{proofId} to the browser as evidence that validation
        succeeded.
\end{itemize}

A separate endpoint allows the receiving system to fetch proofs by identifier.
The service never exposes raw inputs via this endpoint, only proofs and public
signals.

\subsection{Receiving System}

The receiving system, implemented in \texttt{receiving\_system.js}, is the
ultimate arbiter of whether a submission is accepted. It:

\begin{itemize}[leftmargin=*]
  \item accepts a JSON request from the browser containing \texttt{proofId};
  \item fetches corresponding proofs from the validation service via an HTTP
        call to \texttt{/proof/:id};
  \item loads verification keys for each circuit; and
  \item verifies each proof using snarkjs before returning an accept/reject
        decision.
\end{itemize}

By design, the receiving system does not need the raw form data and can make
its decision based solely on the set of verified proofs.

\section{Protocols and Algorithms}

This section describes the main PVE-ZKP workflows in algorithmic form. The
algorithms abstract away implementation details while reflecting the behavior
of the prototype services and client code.

\subsection{Browser Submission Protocol}

\textbf{Protocol 1 (Browser Submission).} The browser client executes the
following steps when the user submits the form:

\begin{enumerate}[leftmargin=*]
  \item Collect form fields \texttt{name}, \texttt{dob}, \texttt{role},
        \texttt{filePath}, and \texttt{token} from the HTML form.
  \item Construct a JSON object \texttt{formData} containing these fields.
  \item Send an HTTP \texttt{POST} request to the validation service at
        \texttt{/validate} on port~3001 with body
        \texttt{"{"formData": formData}"}.
  \item If the response indicates validation failure, display an error message
        to the user and abort.
  \item If the response indicates success and includes a \texttt{proofId},
        construct a second JSON object containing only this identifier.
  \item Send an HTTP \texttt{POST} request to the receiving system at
        \texttt{/receive} on port~3002 with body
        \texttt{"{"proofId": proofId}"}.
  \item Display the receiving system's accept/reject decision to the user.
\end{enumerate}

\subsection{Validation Service Procedure}

\textbf{Protocol 2 (Validation Service).} On receiving a request at
\texttt{/validate}, the validation service performs:

\begin{enumerate}[leftmargin=*]
  \item Parse the JSON body and extract \texttt{formData} fields.
  \item Perform basic well-formedness checks (e.g., required fields present);
        reject if checks fail.
  \item Derive circuit inputs, such as encoding strings as fixed-length ASCII
        arrays and computing age from date of birth.
  \item For each of the five policy circuits (format, threshold, membership,
        integrity, boundary), construct the corresponding input object and call
        the Groth16 proving routine to obtain a proof and public signals.
  \item Aggregate the resulting proofs and signals into a single object and
        store it in an in-memory map keyed by a freshly generated
        \texttt{proofId}.
  \item Return a JSON response indicating success and containing \texttt{proofId}.
\end{enumerate}

\subsection{Receiving System Verification Procedure}

\textbf{Protocol 3 (Receiving System).} On receiving a request at
\texttt{/receive}, the receiving system operates as follows:

\begin{enumerate}[leftmargin=*]
  \item Parse the JSON body and extract \texttt{proofId}; reject if it is
        missing.
  \item Send an HTTP \texttt{GET} request to the validation service's
        \texttt{/proof/:id} endpoint to retrieve the stored proofs and public
        signals associated with \texttt{proofId}.
  \item For each expected policy circuit, load the corresponding verification
        key from the local key store.
  \item For each circuit, invoke the Groth16 verification routine with the
        verification key, public signals, and proof.
  \item If any verification fails or any expected proof is missing, return a
        JSON response indicating that the submission is rejected, along with a
        reason.
  \item If all verifications succeed, return a JSON response indicating that
        the submission is accepted.
\end{enumerate}

These protocols make explicit how the three components cooperate and where
cryptographic assumptions enter the design.

\section{Deployment Considerations}

Although the prototype runs both services on localhost, PVE-ZKP is intended to
support more general deployment patterns. For example, the validation service
could be a third-party microservice operated by a separate organization, while
receiving systems belong to multiple tenants. In such settings, the
architecture's clear separation of roles and data flows becomes particularly
important, and additional concerns such as authentication, authorization, and
rate limiting must be considered.

The current prototype uses an in-memory proof store and a simple HTTP
interface. A production deployment would likely replace the in-memory store
with a durable database or cache (e.g., Redis with appropriate access control
and expiration policies) and add mutual authentication between services. These
engineering extensions do not change the core architecture but strengthen its
operational security.
